\section{Data Transformation}
\label{sec:data_transformation}

\subsection{Forest Mask Application and Quality Filtering}
Following the spatiotemporal ingestion and indexing phase, the raw reflectance data is retrieved from the object storage using the database index. To ensure index values represent true forest dynamics, a sequence of spatial filtering and quality masking is applied dynamically for every query.


% This "late-binding" transformation approach avoids the high storage costs of maintaining pre-processed forest-masked datasets while ensuring that the latest land-cover masks and quality algorithms are always utilized during analysis.

\begin{itemize}
    \item \textbf{Data Source:} Uses ESA WorldCover 10m land cover map.
    \item \textbf{Method:} Apache Sedona spatial join between HLS pixel centroids and WorldCover polygons.
    \item \textbf{Filtering:} Removes non-forest pixels (cropland, grassland, water, built-up surfaces) to reduce computational load.
\end{itemize}

\subsection{Vegetation Index Derivation}
\label{subsec:vi_derivation}
Computes spectral indices from Harmonized Landsat--Sentinel-2 (HLS) data using PySpark.
// formual
\begin{itemize}
    \item \textbf{Indices:} NDVI (Normalized Difference Vegetation Index) and NDMI (Normalized Difference Moisture Index).
    \item \textbf{Bands:} Derived from Red, Near-Infrared (NIR), and Shortwave Infrared (SWIR).
\end{itemize}

\subsection{Temporal Harmonization and Gap Handling}
Standardizes irregular satellite observations into a consistent, regular temporal sequence.
\begin{itemize}
    \item \textbf{Aggregation:} Monthly composites to smooth noise and ensure comparability.
    \item \textbf{Gap Filling:} Short-gap interpolation for missing months.
    \item \textbf{Result:} Produces stable, evenly spaced time series for subsequent analysis.
\end{itemize}

\subsection{Baseline Construction for Anomaly Detection}
Establishes a historical baseline to identify deviations in vegetation behavior.
\begin{itemize}
    \item \textbf{Calculation:} Aggregates monthly NDVI/NDMI values over a fixed 10-year historical period.
    \item \textbf{Metrics:} Calculates the mean and standard deviation for each month.
    \item \textbf{Output:} A climatological signature of normal vegetation behavior for each pixel.
\end{itemize}
% \subsection{Forest Mask Application and Quality Filtering}
% After acquiring multispectral surface reflectance imagery, to ensure that index values represent true forest dynamics, the ESA WorldCover 10\,m land cover map is applied as a forest mask. This is implemented via an Apache Sedona spatial join between the HLS pixel centroids and the WorldCover polygons. By dropping non-forest pixels (cropland, grassland, water, and built-up surfaces) early in the pipeline, the system significantly reduces the computational burden for subsequent time-series analysis. 
% % The literature emphasizes the importance of such masking, noting that NDVI alone ``does not distinguish among vegetation types,'' and that restricting analysis to forest pixels improves interpretability and reduces contamination \cite{worldcover}. Additionally, cloud-affected observations are filtered using Spark DataFrame operations to evaluate scene metadata and quality flags, ensuring that only valid, atmospherically corrected reflectance values contribute to the time series.

% \subsection{Vegetation Index Derivation}
%  The system utilizes PySpark in conjunction with raster libraries (such as Rasterio) to compute Normalized Difference Vegetation Index (NDVI) and the Normalized Difference Moisture Index (NDMI), from the red, near-infrared, and shortwave infrared bands of the Harmonized Landsat--Sentinel-2 (HLS) dataset. 
% % This approach aligns with prior studies showing that long-term vegetation and moisture-sensitive indices can reveal subtle stress signals before visible symptoms emerge (\cite{ndvi}).



% \subsection{Temporal Harmonization and Gap Handling}
% % Because Sentinel--2 and Landsat observations arrive at irregular intervals due to clouds, revisit cycles, and acquisition geometry, the preprocessing pipeline standardizes the data into a regular temporal sequence. The Harmonized Landsat--Sentinel-2 (HLS) product already resolves cross-sensor differences by aligning spectral bands, normalizing reflectance, and placing all imagery on a common grid, so no additional spectral harmonization is required. The pipeline instead aggregates all valid observations within each period---typically a month---into a single composite to smooth noise and ensure year-to-year comparability. When months lack cloud-free data, values are filled through short-gap interpolation or left missing when coverage is insufficient. Producing a stable, evenly spaced NDVI and NDMI time series is necessary for constructing historical baselines and supports statistical models such as BFAST \cite{bfast} and Adaptive EWMA \cite{forests-survey}, which require regular sampling to detect deviations from expected seasonal patterns.
% Because Sentinel--2 and Landsat observations arrive at irregular intervals, the preprocessing pipeline standardizes the data into a regular temporal sequence. The pipeline aggregates valid observations within each period into a single composite to smooth noise and ensure comparability. When months lack cloud-free data, values are filled through short-gap interpolation or left missing when coverage is insufficient. This produces a stable, evenly spaced NDVI and NDMI time series for subsequent analysis.

% \subsection{Baseline Construction for Anomaly Detection}
% For each forest pixel, a multi-year baseline is computed by aggregating monthly NDVI and NDMI values across a fixed historical period (e.g., a 10-year period). The baseline includes the mean and standard deviation for each month, forming a climatological signature of normal vegetation behavior. 
% % This structure mirrors the ``historical seasonal baseline'' approach used in near-real-time anomaly frameworks \cite{nrt-anomalies}, where deviations from expected monthly patterns are converted into standardized anomaly scores. The resulting baselines serve as the foundation for detecting statistically significant vegetation stress preceding disease outbreaks.

\subsection{On demand anomaly detection}
\label{sec:on_demand_detection}

To ensure the system remains responsive and memory-efficient, the transformation and detection logic are not pre-calculated; instead, they are executed on-the-fly for the specific spatiotemporal window requested by the user.

\subsection{Query-Triggered Anomaly detection}
\label{subsec:on_demand_indices}
\subsection{Standardized Z-Score Anomaly Logic}
\label{subsec:z_score_logic}
Immediately following index derivation, the anomaly detection logic is applied to the filtered dataset to identify significant vegetation stress.

\begin{itemize}
    \item \textbf{Statistical Normalization:} For every pixel in the user's query timeframe, the observed index value is compared against a pre-computed historical baseline (2015--2025). The Z-score is calculated as:
    \begin{equation}
        Z = \frac{\text{Observed Value} - \mu_{monthly}}{\sigma_{monthly}}
    \end{equation}
    where $\mu_{monthly}$ and $\sigma_{monthly}$ are the mean and standard deviation for that specific pixel and month.
    \item \textbf{Anomaly Thresholding:} A pixel is flagged as an anomaly if $Z < -2.0$, signaling a statistically rare decline in vegetation health relative to historical seasonal patterns.
    \item \textbf{Query Output:} The system returns the results as either a \textit{Raster Anomaly Map} for spatial visualization or a \textit{Tabular Event Dataset} in Parquet format, providing immediate feedback for the user's specific query.
\end{itemize}